% ============================================================
% Section 81: 有效质量张量的导出与对角化
% ============================================================
\section{固体理论：有效质量张量的导出与主轴方向}
\label{sec:effective-mass-tensor}

\begin{modelbox}[题目：各向异性能带极值点的有效质量张量]
若已知
\[
E(\bm{k})=Ak^2+C(k_xk_y+k_yk_z+k_zk_x),
\]
导出 $\bm{k}=0$ 点的有效质量张量，并找出主轴方向。
\end{modelbox}

\begin{tipbox}[考点快速分析]
\begin{itemize}
    \item \textbf{有效质量逆张量}：$(1/m^*)_{ij}=\hbar^{-2}\partial^2E/\partial k_i\partial k_j|_{\bm{k}=0}$
    \item \textbf{矩阵对角化}：求特征值（有效质量主值）和特征向量（主轴方向）
    \item \textbf{对称矩阵}：$3\times3$ 实对称矩阵，特征值为实数，特征向量正交
    \item \textbf{高对称性矩阵}：所有对角元相等、所有非对角元相等，可快速求特征值
\end{itemize}
\end{tipbox}

\begin{derivationbox}[标准解题过程]
\textbf{Step 1: 有效质量逆张量}

在能带极值点 $\bm{k}=0$ 附近，能量展开到二阶：
\begin{equation}
E(\bm{k})=E(0)+\frac12\sum_{i,j}\Bigl(\frac{\partial^2E}{\partial k_i\partial k_j}\Bigr)_{\!\bm{k}=0}k_ik_j.
\end{equation}

有效质量逆张量
\begin{equation}
\Bigl(\frac{1}{m^*}\Bigr)_{ij}=\frac{1}{\hbar^2}\Bigl(\frac{\partial^2E}{\partial k_i\partial k_j}\Bigr)_{\!\bm{k}=0}.
\end{equation}

对 $E(\bm{k})=A(k_x^2+k_y^2+k_z^2)+C(k_xk_y+k_yk_z+k_zk_x)$ 求二阶偏导：
\begin{equation}
\frac{\partial^2E}{\partial k_x^2}=\frac{\partial^2E}{\partial k_y^2}=\frac{\partial^2E}{\partial k_z^2}=2A,\qquad
\frac{\partial^2E}{\partial k_i\partial k_j}=C\;(i\neq j).
\end{equation}

因此
\begin{equation}
\boxed{\frac{1}{m^*}=\frac{1}{\hbar^2}
\begin{pmatrix}
2A & C & C \\[4pt]
C & 2A & C \\[4pt]
C & C & 2A
\end{pmatrix}.}
\end{equation}

\textbf{Step 2: 特征值与有效质量主值}

解特征方程 $\det\bigl(\frac{1}{m^*}-\lambda I\bigr)=0$。令 $P$ 为上述矩阵，则
\begin{equation}
\det(P-\lambda I)=\begin{vmatrix}
2A-\lambda & C & C \\[4pt]
C & 2A-\lambda & C \\[4pt]
C & C & 2A-\lambda
\end{vmatrix}=0.
\end{equation}

利用矩阵结构（所有对角元相同、所有非对角元相同），特征值可分解为：
\begin{itemize}
    \item 二重根：$\lambda_1=\lambda_2=2A-C$
    \item 单根：$\lambda_3=2A+2C$
\end{itemize}
（验证：迹 $=6A=2(2A-C)+(2A+2C)$；行列式 $=(2A-C)^2(2A+2C)=8A^3-6AC^2+2C^3$，与直接展开一致。）

有效质量主值为特征值的倒数乘以 $\hbar^2$：
\begin{equation}
\boxed{m_1^*=m_2^*=\frac{\hbar^2}{2A-C},\qquad m_3^*=\frac{\hbar^2}{2A+2C}.}
\end{equation}

\textbf{Step 3: 主轴方向（特征向量）}

\textbf{对 $\lambda_3=2A+2C$}：解 $(P-\lambda_3I)\bm{v}=0$
\begin{equation}
\begin{pmatrix}
-2C & C & C \\[4pt]
C & -2C & C \\[4pt]
C & C & -2C
\end{pmatrix}
\begin{pmatrix}x\\y\\z\end{pmatrix}=0
\quad\Longrightarrow\quad x=y=z.
\end{equation}
取单位化特征向量
\begin{equation}
\hat{\bm{e}}_3=\frac{1}{\sqrt{3}}(1,1,1)^\top.
\end{equation}
方向沿立方晶胞的体对角线 $\langle111\rangle$。

\textbf{对 $\lambda_1=\lambda_2=2A-C$}：解 $(P-\lambda_1I)\bm{v}=0$
\begin{equation}
\begin{pmatrix}
C & C & C \\[4pt]
C & C & C \\[4pt]
C & C & C
\end{pmatrix}
\begin{pmatrix}x\\y\\z\end{pmatrix}=0
\quad\Longrightarrow\quad x+y+z=0.
\end{equation}
该平面内任意两个正交向量均可。选取并正交化：
\begin{equation}
\hat{\bm{e}}_1=\frac{1}{\sqrt{2}}(1,0,-1)^\top,\qquad
\hat{\bm{e}}_2=\frac{1}{\sqrt{6}}(1,-2,1)^\top.
\end{equation}
二者均垂直于 $\hat{\bm{e}}_3$，且相互正交（$\hat{\bm{e}}_1\cdot\hat{\bm{e}}_2=0$）。它们张成的平面垂直于 $\langle111\rangle$ 方向。
\end{derivationbox}

\begin{formulabox}[答案汇总]
\begin{center}
\begin{tabular}{ll}
\toprule
\textbf{项目} & \textbf{结果} \\
\midrule
逆有效质量张量 & $\dfrac{1}{\hbar^2}\begin{pmatrix}2A&C&C\\C&2A&C\\C&C&2A\end{pmatrix}$ \\
有效质量主值 & $m_1^*=m_2^*=\dfrac{\hbar^2}{2A-C}$，$m_3^*=\dfrac{\hbar^2}{2A+2C}$ \\
主轴方向 & $\hat{\bm{e}}_3=\dfrac{1}{\sqrt{3}}(1,1,1)$（体对角线） \\
 & $\hat{\bm{e}}_1=\dfrac{1}{\sqrt{2}}(1,0,-1)$，$\hat{\bm{e}}_2=\dfrac{1}{\sqrt{6}}(1,-2,1)$（垂直于体对角线） \\
\bottomrule
\end{tabular}
\end{center}
\end{formulabox}

\begin{insightbox}[物理图像]
\begin{itemize}
    \item 能带极值点附近，等能面为\textbf{椭球面}（三维）或椭圆（二维）。有效质量张量的主轴方向即椭球的主轴方向。
    \item 矩阵 $P$ 的高对称性（所有对角元相等、所有非对角元相等）源于 $E(\bm{k})$ 对 $x,y,z$ 的置换对称性。这种对称性导致一个沿 $\langle111\rangle$ 方向的单重轴和垂直于它的二重简并面。
    \item 若 $C=0$，则 $m_1^*=m_2^*=m_3^*=\hbar^2/(2A)$，各向同性，椭球退化为球面。$C\neq0$ 时的各向异性是晶格非立方对称性或能带高阶耦合的结果。
    \item 有效质量主值的符号决定载流子类型：$m^*>0$ 对应电子（导带底），$m^*<0$ 对应空穴（价带顶）。本题中若 $2A>C>0$，则所有主值为正，$\bm{k}=0$ 为导带底。
\end{itemize}
\end{insightbox}
